% COMP90038 Algorithms and Complexity -- Assignment 1, Semester 2 2026
% Compile with pdfLaTeX (compile out of the box on Overleaf).
\documentclass[11pt,a4paper]{article}
\usepackage[T1]{fontenc}
\usepackage[utf8]{inputenc}
\usepackage{lmodern}
\usepackage[margin=2.3cm,top=2cm,bottom=2.2cm]{geometry}
\usepackage{amsmath,amssymb}
\usepackage{microtype}
\usepackage{enumitem}
\usepackage{float}
\usepackage{algorithm}
\usepackage[noend]{algpseudocode}
% --- house style -------------------------------------------------------------
\algrenewcommand\textproc{\textrm}          % procedure names in roman, as in the spec
\algrenewcommand\algorithmicrequire{\textbf{input:}}
\renewcommand{\thealgorithm}{\arabic{algorithm}}
\newcommand{\id}[1]{\ensuremath{\mathit{#1}}}   % identifiers
\newcommand{\Oh}{\ensuremath{O}}
\setlength{\parindent}{0pt}
\setlength{\parskip}{0.45em}
\setlist{topsep=0.2em,itemsep=0.15em,parsep=0pt}
\setlength{\intextsep}{0.7\baselineskip}   % space around the algorithm floats
\begin{document}
% --- title -------------------------------------------------------------------
\begin{center}
Name: JIAPING BAI \qquad Student ID: 1864707
\end{center}
\vspace{-0.4em}
\hrule
\vspace{0.6em}
%==============================================================================
\section*{Problem 1}
\subsection*{(a)(i)}
\[
T(0)=\Theta(1),
\qquad
T(n)=T(n-1)+\Theta(\log n) \quad\text{for } n\ge 1 .
\]
\subsection*{(a)(ii)}
$T(n)\in\Theta(n\log n)$. Backward substitution gives
$T(n)=T(0)+\sum_{i=1}^{n}\Theta(\log i)=\Theta\bigl(\log_2(n!)\bigr)$, using
$\sum_{i=1}^{n}\log_2 i=\log_2\prod_{i=1}^{n} i=\log_2(n!)$.

By Stirling's formula, $n!=\sqrt{2\pi n}\,(n/e)^{n}\bigl(1+\Theta(1/n)\bigr)$, so
$\log_2(n!)=n\log_2 n-n\log_2 e+\Oh(\log n)=n\log_2 n-\Theta(n)$.
The subtracted term is positive, so $\log_2(n!)\le n\log_2 n$ and $\log_2(n!)\in\Oh(n\log n)$; and
since $\Theta(n)=o(n\log n)$, we get $\log_2(n!)\ge n\log_2 n-o(n\log_2 n)$, so
$\log_2(n!)\in\Omega(n\log n)$. Hence $\log_2(n!)\in\Theta(n\log n)$ and therefore
$T(n)\in\Theta(n\log n)$.
\subsection*{(b)(i)}
\textbf{return } $2p+q,\;2q$
\subsection*{(b)(ii)}
$T(k)\in\Theta(k)$. Each invocation of \textit{TotalLightCountBitsRec} does constant work and makes
one recursive call on $k-1$, so $T(k)=T(k-1)+\Theta(1)$ with $T(1)=\Theta(1)$, which unwinds to
$\Theta(k)$. \textit{TotalLightCountBits} adds one call and constant work, so it is $\Theta(k)$ as
well.
\subsection*{(b)(iii)}
\textit{TotalLightCountBits} is more efficient, by a factor of $\Theta(n)$: $\Theta(\log n)$
against $\Theta(n\log n)$.

Let $n=2^{k}$, so $k=\log_2 n$. The integers representable with $k$ bits are exactly
$0,\dots,2^{k}-1=0,\dots,n-1$, i.e.\ all the tram numbers except $n$ itself; and $n=2^{k}$ has
exactly one light on. So the answer to the problem of part~(a) is
\[
\textit{TotalLightCountBits}(\log_2 n) + 1 .
\]
Its cost is $\Theta(\log n)$ to obtain the binary representation of $n$, plus
$\Theta(k)=\Theta(\log n)$ for the recursion by (b)(ii): in total $\Theta(\log n)$.
%==============================================================================
\section*{Problem 2}
\subsection*{(a)(i)}
\[
\id{arr}=[1,\,1]
\qquad\text{and}\qquad
\id{start}=0 .
\]

Since no entry of \id{arr} is $0$, the network has no terminal stop, so the configuration is
invalid and the procedure must return \textbf{False}. \textit{IsValid\_AI} instead recurses
forever, so it never returns \textbf{False}.

\subsection*{(a)(ii)}

The disjunction \textbf{or} is evaluated from left to right, so the call
\textit{IsValid\_AI}$(\id{arr},\,i+\id{dist})$ is made first.
\begin{enumerate}[label=\arabic*.,leftmargin=2.2em]
\item \textit{IsValid\_AI}$([1,1],\,0)$: $i=0$; $0\le i<2$; $\id{arr}[0]=1\neq 0$;
$\id{dist}\gets 1$; evaluate the first disjunct, i.e.\ call \dots
\item \textit{IsValid\_AI}$([1,1],\,1)$: $i=1$; $0\le i<2$; $\id{arr}[1]=1\neq 0$;
$\id{dist}\gets 1$; evaluate the first disjunct, i.e.\ call \dots
\item \textit{IsValid\_AI}$([1,1],\,2)$: $i=2\ge \id{length}(\id{arr})=2$, so
\textbf{return False}.
\item Back in call~2: the first disjunct is \textbf{False}, so evaluate the second one, i.e.\
call \textit{IsValid\_AI}$([1,1],\,0)$ --- and we are back in call~1.
\item \textit{IsValid\_AI}$([1,1],\,0)$: identical to call~1; evaluate the first disjunct \dots
\item \textit{IsValid\_AI}$([1,1],\,1)$: identical to call~2; evaluate the first disjunct \dots
\item \textit{IsValid\_AI}$([1,1],\,2)$: \textbf{return False}.
\item Back in call~6: evaluate the second disjunct, i.e.\ call
\textit{IsValid\_AI}$([1,1],\,0)$ --- and the whole pattern starts again.
\end{enumerate}
Calls 5--8 repeat calls 1--4 with exactly the same arguments, so the sequence of calls
$0,1,2,0,1,2,\dots$ continues forever and the top-level call never returns.

\subsection*{(b)}

A depth-first search from \id{start}, marking the stops already visited.

\begin{algorithm}[H]
\caption{}
\begin{algorithmic}[1]
\Function{IsValid}{$\id{arr},\, \id{start}$}
\State $n \gets \Call{length}{\id{arr}}$
\If{$\id{start} < 0$ \textbf{or} $\id{start} \ge n$}
\State \Return False
\EndIf
\State $\id{visited} \gets$ \textbf{new} Boolean array with indices $0 \ldots n-1$
\For{$i \gets 0$ \textbf{to} $n-1$}
\State $\id{visited}[i] \gets$ False
\EndFor
\State \Return $\Call{Search}{\id{arr},\, \id{start},\, \id{visited}}$
\EndFunction
\Statex
\Function{Search}{$\id{arr},\, i,\, \id{visited}$} \Comment{\id{visited} is shared by all calls}
\If{$i < 0$ \textbf{or} $i \ge \Call{length}{\id{arr}}$}
\State \Return False
\EndIf
\If{$\id{visited}[i]$}
\State \Return False
\EndIf
\State $\id{visited}[i] \gets$ True
\If{$\id{arr}[i] = 0$}
\State \Return True
\EndIf
\State $\id{dist} \gets \id{arr}[i]$
\If{$\Call{Search}{\id{arr},\, i + \id{dist},\, \id{visited}}$}
\State \Return True
\EndIf
\State \Return $\Call{Search}{\id{arr},\, i - \id{dist},\, \id{visited}}$
\EndFunction
\end{algorithmic}
\end{algorithm}
\subsection*{(c)}
Since $\id{visited}[i]$ is set before the recursive calls and never reset, at most $n$ calls to
\textit{Search} get past the visited test. Each of those does $\Oh(1)$ work and makes at most two
more calls, so in total at most $2n+1$ calls are made, each taking $\Oh(1)$. Adding the
$\Theta(n)$ initialisation of \id{visited} gives $\Oh(n)$.
%==============================================================================
\section*{Problem 3}
\subsection*{(a)}
Within $t$ minutes tram $i$ makes $\lfloor t/\id{time}[i]\rfloor$ trips.
\begin{algorithm}[H]
\caption{}
\begin{algorithmic}[1]
\Function{MaxTrips}{$\id{time},\, t$}
\State $n \gets \Call{length}{\id{time}}$
\State $\id{trips} \gets 0$
\For{$i \gets 0$ \textbf{to} $n-1$}
\State $\id{trips} \gets \id{trips} + \lfloor t / \id{time}[i] \rfloor$
\EndFor
\State \Return $\id{trips}$
\EndFunction
\end{algorithmic}
\end{algorithm}
$\textit{MaxTrips}(\id{time},\cdot)$ is \emph{non-decreasing} in $t$, since every
$\lfloor t/\id{time}[i]\rfloor$ is.
\subsection*{(b)}
\begin{algorithm}[H]
\caption{}
\begin{algorithmic}[1]
\Function{MinTimeBruteForce}{$\id{time},\, m$}
\State $t \gets 1$
\While{$\Call{MaxTrips}{\id{time}, t} < m$}
\State $t \gets t + 1$
\EndWhile
\State \Return $t$
\EndFunction
\end{algorithmic}
\end{algorithm}
The loop terminates: writing $t_{\min}=\min_i \id{time}[i]$, the fastest tram on its own makes
$m$ trips within $t_{\min}\cdot m$ minutes, so $\textit{MaxTrips}(\id{time}, t_{\min}m)\ge m$ and
the loop stops at some $t\le t_{\min}\cdot m$.
\subsection*{(c)}
Let $T$ be the value returned. The loop condition is evaluated once for each of
$t=1,2,\dots,T$, and each evaluation runs \textit{MaxTrips}, which costs
$\Theta(n)$; the rest of each iteration is $\Oh(1)$. So the running time is $\Theta(n\cdot T)$.

Since $\id{time}[i]\le k$ for every $i$,
\begin{align*}
\textit{MaxTrips}\bigl(\id{time},&k\lceil m/n\rceil\bigr)  =\sum_{i}\left\lfloor \frac{k\lceil m/n\rceil}{\id{time}[i]}\right\rfloor \\
\;&\ge\; n\lceil m/n\rceil \;\ge\; m ,
\end{align*}
so $T\le k\lceil m/n\rceil$; and this bound is attained when every tram is as slow as possible,
$\id{time}[i]=k$ for all $i$, for then $\textit{MaxTrips}(\id{time},t)=n\lfloor t/k\rfloor$, which
first reaches $m$ at $t=k\lceil m/n\rceil$. Hence the worst-case running time is
\[
\Theta\bigl(n\cdot k\lceil m/n\rceil\bigr)=\Theta\bigl(k(m+n)\bigr), \qquad\text{since } m\le n\lceil m/n\rceil < m+n .
\]
\subsection*{(d)}
By (a) the values $\textit{MaxTrips}(\id{time},1),\textit{MaxTrips}(\id{time},2),\dots$ form a
sorted array, and $T$ is the position of the first entry that reaches $m$. Binary search finds
that position. Part~(b) supplies the right end of the range, $t_{\min}\cdot m$.
\begin{algorithm}[H]
\caption{}
\begin{algorithmic}[1]
\Function{MinTimeBinarySearch}{$\id{time},\, m$}
\State $n \gets \Call{length}{\id{time}}$
\State $t_{\min} \gets \id{time}[0]$
\For{$i \gets 1$ \textbf{to} $n-1$}
\If{$\id{time}[i] < t_{\min}$}
\State $t_{\min} \gets \id{time}[i]$
\EndIf
\EndFor
\State $\id{lo} \gets 1$
\State $\id{hi} \gets t_{\min} \cdot m$
\Comment{$\textit{MaxTrips}(\id{time},\id{hi})\ge m$, so $T\in[\id{lo},\id{hi}]$}
\While{$\id{lo} < \id{hi}$}
\State $\id{mid} \gets \lfloor (\id{lo}+\id{hi})/2 \rfloor$
\If{$\Call{MaxTrips}{\id{time}, \id{mid}} \ge m$}
\State $\id{hi} \gets \id{mid}$
\Else
\State $\id{lo} \gets \id{mid} + 1$
\EndIf
\EndWhile
\State \Return $\id{lo}$
\EndFunction
\end{algorithmic}
\end{algorithm}
\subsection*{(e)}
Since computing $t_{\min}$ takes $\Theta(n)$, and the search interval initially contains
$\id{hi}-\id{lo}+1=t_{\min}\cdot m\le k\,m$ candidate times (recall $t_{\min}\le k$), and every
iteration at least halves it, there are at most $\lceil\log_2(km)\rceil+1$ iterations. Each
iteration is dominated by its single call to \textit{MaxTrips}, i.e.\ $\Theta(n)$; the rest is
$\Oh(1)$. Hence the worst-case running time is
\[
\Theta(n)+\Oh\bigl(n\log(km)\bigr)=\Oh\bigl(n\log(km)\bigr)=\Oh\bigl(n(\log k+\log m)\bigr).
\]
\end{document}
